\chapter{工具系统 (Tool System)}

\section{架构概述}

工具系统是 Claude Code 的核心特性——它赋予了 LLM "手脚"，使其能够真正执行软件工程任务。所有工具的注册和管理通过 \code{src/tools.ts} 完成，而工具的类型定义在 \code{src/Tool.ts}（约 29,000 行）中。

\section{工具类型定义}

\subsection{核心接口}

每个工具必须实现一组标准接口：

\begin{lstlisting}[language=TypeScript,breaklines=true]
export type ToolInputJSONSchema = {
  type: 'object'
  properties?: { [x: string]: unknown }
  [x: string]: unknown
}

// Tool 的核心字段（从源码推断）
interface Tool {
  name: string                        // 工具名称（唯一标识）
  description: string                 // 工具描述（给 LLM 看的）
  inputSchema: ToolInputJSONSchema    // 输入参数 JSON Schema
  isEnabled(): boolean                // 是否启用
  
  // 权限相关
  mcpInfo?: {                         // MCP 工具信息
    serverName: string
    toolName: string
  }
  
  // 执行相关
  maxResultSizeChars?: number         // 最大结果大小
}
\end{lstlisting}

\subsection{权限上下文}

\begin{lstlisting}[language=TypeScript,breaklines=true]
export type ToolPermissionContext = DeepImmutable<{
  mode: PermissionMode
  additionalWorkingDirectories: Map<string, AdditionalWorkingDirectory>
  alwaysAllowRules: ToolPermissionRulesBySource
  alwaysDenyRules: ToolPermissionRulesBySource
  alwaysAskRules: ToolPermissionRulesBySource
  isBypassPermissionsModeAvailable: boolean
  isAutoModeAvailable?: boolean
  strippedDangerousRules?: ToolPermissionRulesBySource
  shouldAvoidPermissionPrompts?: boolean
  awaitAutomatedChecksBeforeDialog?: boolean
  prePlanMode?: PermissionMode
}>
\end{lstlisting}

关键字段说明：

\begin{longtable}{ll}
\toprule
字段 & 说明 \\
\midrule
\code{mode} & 权限模式：\code{default} / \code{plan} / \code{bypassPermissions} / \code{auto} \\
\code{alwaysAllowRules} & 始终允许的规则（按来源分组） \\
\code{alwaysDenyRules} & 始终拒绝的规则（按来源分组） \\
\code{shouldAvoidPermissionPrompts} & 后台代理使用：无法显示 UI 时自动拒绝 \\
\code{prePlanMode} & 进入计划模式前的权限模式，退出时恢复 \\
\bottomrule
\end{longtable}

\section{工具注册表 (tools.ts)}

\subsection{静态工具注册}

\begin{lstlisting}[language=TypeScript,breaklines=true]
import { AgentTool } from './tools/AgentTool/AgentTool.js'
import { BashTool } from './tools/BashTool/BashTool.js'
import { FileEditTool } from './tools/FileEditTool/FileEditTool.js'
import { FileReadTool } from './tools/FileReadTool/FileReadTool.js'
import { FileWriteTool } from './tools/FileWriteTool/FileWriteTool.js'
import { GlobTool } from './tools/GlobTool/GlobTool.js'
import { GrepTool } from './tools/GrepTool/GrepTool.js'
import { WebFetchTool } from './tools/WebFetchTool/WebFetchTool.js'
import { WebSearchTool } from './tools/WebSearchTool/WebSearchTool.js'
import { SkillTool } from './tools/SkillTool/SkillTool.js'
// ... 更多工具
\end{lstlisting}

\subsection{条件工具注册}

许多工具通过 Feature Flag 或环境变量条件加载：

\begin{lstlisting}[language=TypeScript,breaklines=true]
// 仅内部用户可用
const REPLTool = process.env.USER_TYPE === 'ant'
  ? require('./tools/REPLTool/REPLTool.js').REPLTool
  : null

// 需要 Feature Flag
const SleepTool = feature('PROACTIVE') || feature('KAIROS')
  ? require('./tools/SleepTool/SleepTool.js').SleepTool
  : null

// 触发器相关工具
const cronTools = feature('AGENT_TRIGGERS')
  ? [
      require('./tools/ScheduleCronTool/CronCreateTool.js').CronCreateTool,
      require('./tools/ScheduleCronTool/CronDeleteTool.js').CronDeleteTool,
      require('./tools/ScheduleCronTool/CronListTool.js').CronListTool,
    ]
  : []

// 协调器模式
const coordinatorModeModule = feature('COORDINATOR_MODE')
  ? require('./coordinator/coordinatorMode.js')
  : null

// Web 浏览器工具
const WebBrowserTool = feature('WEB_BROWSER_TOOL')
  ? require('./tools/WebBrowserTool/WebBrowserTool.js').WebBrowserTool
  : null
\end{lstlisting}

\subsection{懒加载（打破循环依赖）}

\begin{lstlisting}[language=TypeScript,breaklines=true]
// 懒加载以打破循环依赖：tools.ts -> TeamCreateTool -> ... -> tools.ts
const getTeamCreateTool = () =>
  require('./tools/TeamCreateTool/TeamCreateTool.js').TeamCreateTool
const getTeamDeleteTool = () =>
  require('./tools/TeamDeleteTool/TeamDeleteTool.js').TeamDeleteTool
const getSendMessageTool = () =>
  require('./tools/SendMessageTool/SendMessageTool.js').SendMessageTool
\end{lstlisting}

\subsection{完整工具列表}

\begin{lstlisting}[language=TypeScript,breaklines=true]
export function getAllBaseTools(): Tools {
  return [
    AgentTool,                     // 子代理生成
    TaskOutputTool,                // 任务输出
    BashTool,                      // Shell 命令
    // 嵌入式搜索工具检测
    ...(hasEmbeddedSearchTools() ? [] : [GlobTool, GrepTool]),
    ExitPlanModeV2Tool,            // 退出计划模式
    FileReadTool,                  // 文件读取
    FileEditTool,                  // 文件编辑
    FileWriteTool,                 // 文件写入
    NotebookEditTool,              // Notebook 编辑
    WebFetchTool,                  // URL 获取
    TodoWriteTool,                 // Todo 写入
    WebSearchTool,                 // 网页搜索
    TaskStopTool,                  // 任务停止
    AskUserQuestionTool,           // 向用户提问
    SkillTool,                     // 技能执行
    EnterPlanModeTool,             // 进入计划模式
    // 条件工具...
    ...(isTodoV2Enabled()
      ? [TaskCreateTool, TaskGetTool, TaskUpdateTool, TaskListTool]
      : []),
    ...(isWorktreeModeEnabled()
      ? [EnterWorktreeTool, ExitWorktreeTool]
      : []),
    getSendMessageTool(),          // 代理间消息
    ...(isAgentSwarmsEnabled()
      ? [getTeamCreateTool(), getTeamDeleteTool()]
      : []),
    BriefTool,                     // 简要输出
    ListMcpResourcesTool,          // 列出 MCP 资源
    ReadMcpResourceTool,           // 读取 MCP 资源
    ...(isToolSearchEnabledOptimistic() ? [ToolSearchTool] : []),
    // ... 更多条件工具
  ]
}
\end{lstlisting}

\textbf{嵌入式搜索工具}：在 Anthropic 内部构建中，\code{bfs}（fast find）和 \code{ugrep}（fast grep）被嵌入到 Bun 二进制文件中。当这些工具可用时，\code{GlobTool} 和 \code{GrepTool} 就不需要了——Shell 中的 find/grep 已被别名为这些快速工具。

\section{工具分类详解}

\subsection{文件操作工具}

\begin{longtable}{lll}
\toprule
工具 & 功能 & 关键特性 \\
\midrule
\code{FileReadTool} & 读取文件 & 支持图片、PDF、Notebook；LRU 缓存 \\
\code{FileWriteTool} & 创建/覆盖文件 & 完整文件写入 \\
\code{FileEditTool} & 部分编辑 & 字符串替换，精确匹配 \\
\code{NotebookEditTool} & Jupyter Notebook & 单元格级别操作 \\
\code{GlobTool} & 文件搜索 & picomatch 模式匹配 \\
\code{GrepTool} & 内容搜索 & ripgrep 后端 \\
\bottomrule
\end{longtable}

\subsection{执行工具}

\begin{longtable}{lll}
\toprule
工具 & 功能 & 关键特性 \\
\midrule
\code{BashTool} & Shell 命令 & 支持 Bash 和 PowerShell \\
\code{PowerShellTool} & PowerShell 命令 & Windows 专用 \\
\code{LSPTool} & 语言服务器 & 代码智能：跳转、补全、诊断 \\
\bottomrule
\end{longtable}

\subsection{网络工具}

\begin{longtable}{ll}
\toprule
工具 & 功能 \\
\midrule
\code{WebFetchTool} & 获取 URL 内容 \\
\code{WebSearchTool} & 执行网页搜索 \\
\code{WebBrowserTool} & 浏览器自动化（Feature Flag 控制） \\
\bottomrule
\end{longtable}

\subsection{代理与协调工具}

\begin{longtable}{ll}
\toprule
工具 & 功能 \\
\midrule
\code{AgentTool} & 生成子代理执行复杂任务 \\
\code{TeamCreateTool} & 创建团队级并行代理 \\
\code{TeamDeleteTool} & 删除团队代理 \\
\code{SendMessageTool} & 代理间消息传递 \\
\code{TaskCreateTool} & 创建任务 \\
\code{TaskUpdateTool} & 更新任务状态 \\
\code{TaskStopTool} & 停止运行中的任务 \\
\bottomrule
\end{longtable}

\subsection{模式控制工具}

\begin{longtable}{ll}
\toprule
工具 & 功能 \\
\midrule
\code{EnterPlanModeTool} & 进入计划模式（只读，不执行） \\
\code{ExitPlanModeV2Tool} & 退出计划模式 \\
\code{EnterWorktreeTool} & 进入 Git Worktree 隔离环境 \\
\code{ExitWorktreeTool} & 退出 Worktree \\
\bottomrule
\end{longtable}

\subsection{特殊工具}

\begin{longtable}{ll}
\toprule
工具 & 功能 \\
\midrule
\code{SkillTool} & 执行预定义或自定义技能 \\
\code{ToolSearchTool} & 延迟工具发现（当工具数量多时） \\
\code{SleepTool} & 主动模式下的等待 \\
\code{CronCreateTool} & 创建定时触发器 \\
\code{SyntheticOutputTool} & 结构化输出生成 \\
\code{BriefTool} & 简要输出模式 \\
\code{AskUserQuestionTool} & 向用户提出交互式问题 \\
\bottomrule
\end{longtable}

\section{工具过滤机制}

\subsection{拒绝规则过滤}

\begin{lstlisting}[language=TypeScript,breaklines=true]
export function filterToolsByDenyRules<T extends { name: string; mcpInfo?: ... }>(
  tools: readonly T[],
  permissionContext: ToolPermissionContext,
): T[] {
  return tools.filter(tool => !getDenyRuleForTool(permissionContext, tool))
}
\end{lstlisting}

工具在到达 LLM 之前就被过滤。一个被全面拒绝（blanket deny）的工具根本不会出现在模型的工具列表中。

\subsection{REPL 模式工具隐藏}

\begin{lstlisting}[language=TypeScript,breaklines=true]
if (isReplModeEnabled()) {
  const replEnabled = allowedTools.some(tool =>
    toolMatchesName(tool, REPL_TOOL_NAME)
  )
  if (replEnabled) {
    // 隐藏原始工具——它们通过 REPL 的 VM 上下文间接可用
    allowedTools = allowedTools.filter(
      tool => !REPL_ONLY_TOOLS.has(tool.name)
    )
  }
}
\end{lstlisting}

在 REPL 模式下，底层工具（Bash、FileRead 等）被隐藏，因为它们通过 REPL 的 VM 上下文间接可用。

\subsection{简单模式}

\begin{lstlisting}[language=TypeScript,breaklines=true]
if (isEnvTruthy(process.env.CLAUDE_CODE_SIMPLE)) {
  const simpleTools: Tool[] = [BashTool, FileReadTool, FileEditTool]
  // 协调器模式下额外添加 AgentTool 和 TaskStopTool
  if (feature('COORDINATOR_MODE') && coordinatorModeModule?.isCoordinatorMode()) {
    simpleTools.push(AgentTool, TaskStopTool, getSendMessageTool())
  }
  return filterToolsByDenyRules(simpleTools, permissionContext)
}
\end{lstlisting}

简单模式只提供三个基本工具：Bash、读取、编辑。

\section{工具池组装}

\begin{lstlisting}[language=TypeScript,breaklines=true]
export function assembleToolPool(
  permissionContext: ToolPermissionContext,
  mcpTools: Tools,
): Tools {
  const builtInTools = getTools(permissionContext)
  const allowedMcpTools = filterToolsByDenyRules(mcpTools, permissionContext)

  // 排序以保证 prompt cache 稳定性
  const byName = (a: Tool, b: Tool) => a.name.localeCompare(b.name)
  return uniqBy(
    [...builtInTools].sort(byName).concat(allowedMcpTools.sort(byName)),
    'name',  // 内置工具在名称冲突时优先
  )
}
\end{lstlisting}

组装策略：
\begin{enumerate}[nosep]
  \item 获取内置工具（经过模式过滤和启用检查）
  \item 获取 MCP 工具（经过拒绝规则过滤）
  \item 内置工具作为连续前缀排序
  \item MCP 工具在后面排序
  \item 通过 \code{uniqBy} 去重（内置工具优先级更高）
\end{enumerate}

\section{ToolUseContext 深入}

每次工具执行时传入的上下文对象包含丰富的运行时信息：

\begin{lstlisting}[language=TypeScript,breaklines=true]
export type ToolUseContext = {
  // 配置
  options: {
    commands: Command[]
    tools: Tools
    mcpClients: MCPServerConnection[]
    mcpResources: Record<string, ServerResource[]>
    mainLoopModel: string
    thinkingConfig: ThinkingConfig
    agentDefinitions: AgentDefinitionsResult
    refreshTools?: () => Tools  // MCP 服务器中途连接时刷新
  }
  
  // 控制
  abortController: AbortController      // 中止信号
  readFileState: FileStateCache          // 文件读取缓存
  
  // 状态
  getAppState(): AppState
  setAppState(f: (prev: AppState) => AppState): void
  messages: Message[]                    // 当前对话消息

  // 限制
  fileReadingLimits?: { maxTokens?; maxSizeBytes? }
  globLimits?: { maxResults? }
  
  // 权限追踪
  toolDecisions?: Map<string, { source; decision; timestamp }>
  localDenialTracking?: DenialTrackingState
  
  // UI 交互
  setToolJSX?: SetToolJSXFn             // 设置工具 UI
  addNotification?: (notif) => void     // 添加通知
  sendOSNotification?: (opts) => void   // OS 级别通知
  
  // 子代理
  agentId?: AgentId                      // 子代理 ID
  agentType?: string                     // 代理类型名称
}
\end{lstlisting}

\section{设计启示}

\subsection{声明式工具定义}

每个工具是自包含的模块——定义自己的输入 Schema、权限模型和执行逻辑。这种模式使得添加新工具极其简单，不需要修改任何中心化的路由逻辑。

\subsection{条件加载与死代码消除}

通过 \code{feature()} 和 \code{process.env} 的条件加载，确保未使用的工具代码在构建时被完全移除。这对控制 bundle 大小至关重要。

\subsection{工具池的缓存感知排序}

将内置工具排在前面并保持稳定排序，使得 Prompt Cache 可以跨请求复用。这种"缓存感知"的设计在高频调用场景下可以显著降低成本。

\bigskip\noindent\rule{\textwidth}{0.4pt}\bigskip

\textbf{下一章}：\href{./chapter-07-command-system.md}{命令系统 (Command System)}
